\section{Introduction}
\subsection{Barn owls}
Barn owls (genus \textit{Tyto}) are a group that is found all over the world. This work focuses on the Western barn owl (\textit{Tyto alba}). The population of \textit{T. alba} in Britain is in gradual decline~\cite{keeneStudySmallMammal2009,loveChangesFoodBritish2000}, stressing the need for conservation to be stronger and stuff. 

\subsection{Vocal individuality}
Vocal individuality (VI) in birds has been studied to a significant extent in the literature -- several bird species have been observed to identify mates and siblings via sound~\cite{lindFemaleGreatTits1996,carlsonIndividualVocalRecognition2020}. Owls are no exception, and generally mate for life, necessitating some means for individuals to recognise each other long term. Indeed, owl vocalisations have been shown to be both consistent across time (i.e., the individual could be reliably identified by its vocal characteristics over a period of years), as in the eagle owl (\textit{Bubo bubo})~\cite{lengagneTEMPORALSTABILITYINDIVIDUAL2001}, the African wood owl (\textit{Strix woodfordii})~\cite{delportVocalIdentificationIndividual2002} and the Sunda Scops-Owl (\textit{Otus lempiji})~\cite{yeeVocalIndividualitySunda2016}. There is also evidence that vocalisations are distinct between individuals, as in the Pygmy owl (\textit{Glaucidium passerinum})~\cite{galeottiIndividuallyDistinctHooting1993} and the Eastern screech owl (\textit{Megascops asio})~\cite{nagyIdentificationIndividualEastern2012}. This makes it possible to consider VI-based methods in conservation models to glean information  about the life history and migration of birds -- such information can be critical for making decisions relating to the management of these birds~\cite{terryRoleVocalIndividuality2005,choiCaseStudyMale2019}.

Some of the earliest work aiming to use VI for individual owl classification was carried out by Galeotti et al. on Pygmy owls~\cite{galeottiIndividuallyDistinctHooting1993}, which was used to show that male \textit{G. passerinum} individuals could be reliably distinguished by their territorial hooting calls. These individuals were labelled and therefore classifier predictions could be validated against this ground truth. Turning to \textit{T. alba}, Dreiss et al. demonstrate that barn owl nestlings have evolved identifiable vocal characteristics as a result of sibling competition~\cite{dreissIndividualVocalSignatures2014}. Given this theoretical justification, it is reasonable to assume that it should be possible for barn owl chicks' calls to be differentiated, supporting the viability of a machine learning approach to vocal individuality. Machine learning is particularly well-suited to this problem due to the imperceptible variations in the barn owl vocalisations -- juvenile screech vocalisations are far less distinctive than adult calls, and are qualitatively indistinguishable by both ear and through visual inspection of spectrograms (See Figure~\ref{fig:dsp_pipeline}). 

\subsection{Modern VI approaches}
Machine learning has been a part of bird call analysis for quite some time -- early works employed discriminant analysis for classifying labelled calls~\cite{galeottiIndividuallyDistinctHooting1993,lengagneTEMPORALSTABILITYINDIVIDUAL2001,delportVocalIdentificationIndividual2002,gravaIndividualAcousticMonitoring2008,pamelaIDENTIFICATIONINDIVIDUALBARRED2008,yeeVocalIndividualitySunda2016}, but these approaches are limited because the number of individuals in an ``in the wild'' conservation setting are unknown~\cite{nagyIdentificationIndividualEastern2012}. Moreover, they rely on a few hand-crafted features such as the timing or selected attributes of the vocalisation. As explained earlier, this is not always possible due to the imperceptible nature of juvenile hissing calls in \textit{T. alba}. More recent works have introduced deep learning models, incorporating modern innovations such as attention, deep convolutional networks and variational auto-encoders (VAEs)~\cite{ntalampirasAcousticDetectionUnknown2021,xieHighAccuracyIndividual2020,bedoyaAcousticCensusingIndividual2021} to address these limitations. The work of Stowell et al.~\cite{stowellAutomaticAcousticIdentification2018} is notable for its data augmentation strategy. Given recordings of songs from multiple individuals across several species and locations, the birdsong is first separated from the environmental background noise present at the recording site. Synthetic training samples are then generated by combining the isolated birdsong ``foreground'' with randomly selected background recordings. This approach reduces the likelihood that the model will rely on location-specific environmental noise when learning to discriminate between individuals. All the works discussed above focus on adult individuals and employ ground-truth identity labels when training models. 

Contrary to these applications, our work investigates the VI of barn owl chicks, with the objective of discerning the number of individuals present in the nest. The lack of video equipment within the nest meant that the calls could not be mapped to the individual that made them, necessitating the use of an unsupervised approach.

Vocal individuality in Crested Ibis (\textit{Nipponia nippon})using an autoencoder and attention mechanism~\cite{xieHighAccuracyIndividual2020}. This one uses an RNN based auto-encoder with attention as the encoder, and the latent representation is decoded by a symmetrical decoder, which reconstructs the original signal. A similar approach is used by Ntalampiras et al., who use an encoder-decoder architecture to reconstruct spectrograms from several avian species~\cite{ntalampirasAcousticDetectionUnknown2021}. However, given that all approaches so far are interested in the clustering or classification of calls rather than generating synthetic ones, it makes more sense to use VAEs to reconstruct the latent representations, not all the pixels in the original signal. This has shown to be more effective by Bardes et al. in the field of video classification~\cite{bardesRevisitingFeaturePrediction2024}. We take the latter approach, closely following the strategy of predicting latent blocks as in~\cite{tongVideoMAEMaskedAutoencoders2022}. 

Another study done on birds (not owls) and uses Deep CNN architecture to do vocal individuality~\cite{bedoyaAcousticCensusingIndividual2021}.  Unclear what kind of birds they used. Uses labels/supervised \\

The data collected in this study lacks ground-truth identity labels, as there was no independent mechanism for associating calls with individual owlets. This necessitated the use of a fully unsupervised deep learning approach. \\

- We hypothesise that it is possible to differentiate between different individuals on the basis of vocal characteristics on a quantitative basis, even though they sound qualitatively indistinguishable. 